% Compilação AVULSA de um fragmento (volN/.../<stem>.tex) — p. ex. pelo botão do editor,
% de dentro da pasta do nó.  O fragmento começa com
%
%   \ifdefined\pg\else\IfFileExists{tex/avulso.tex}{% Compilação AVULSA de um fragmento (volN/.../<stem>.tex) — p. ex. pelo botão do editor,
% de dentro da pasta do nó.  O fragmento começa com
%
%   \ifdefined\pg\else\IfFileExists{tex/avulso.tex}{% Compilação AVULSA de um fragmento (volN/.../<stem>.tex) — p. ex. pelo botão do editor,
% de dentro da pasta do nó.  O fragmento começa com
%
%   \ifdefined\pg\else\IfFileExists{tex/avulso.tex}{% Compilação AVULSA de um fragmento (volN/.../<stem>.tex) — p. ex. pelo botão do editor,
% de dentro da pasta do nó.  O fragmento começa com
%
%   \ifdefined\pg\else\IfFileExists{tex/avulso.tex}{\input{tex/avulso}}{\input{../../tex/avulso}}\fi
%
% e termina com  \csname fimavulso\endcsname .  Dentro do livro ou do _preview, \pg já existe
% e nada disto roda; \fimavulso fica indefinido e o \csname vira \relax.
%
% Motor: o projeto é lualatex (fontspec/polyglossia).  Se o editor chamar pdflatex, cai
% num preâmbulo equivalente com inputenc/babel — serve para conferir, não para o PDF final.
\documentclass[11pt,a5paper]{book}
\ifdefined\directlua
  \usepackage{fontspec}
  \usepackage{polyglossia}
  \setmainlanguage{portuguese}
\else\ifdefined\XeTeXrevision
  \usepackage{fontspec}
  \usepackage{polyglossia}
  \setmainlanguage{portuguese}
\else
  \usepackage[utf8]{inputenc}
  \usepackage[T1]{fontenc}
  \usepackage{lmodern}
  \usepackage[portuguese]{babel}
  \DeclareUnicodeCharacter{2192}{\ensuremath{\rightarrow}}
\fi\fi
% raiz do repo: compilado da raiz ou de volN/<parte>/ (dois níveis abaixo)
\IfFileExists{tex/preambulo.tex}{\input{tex/preambulo}}{\input{../../tex/preambulo}}
\def\fimavulso{\end{document}}
\begin{document}
}{\input{../../tex/avulso}}\fi
%
% e termina com  \csname fimavulso\endcsname .  Dentro do livro ou do _preview, \pg já existe
% e nada disto roda; \fimavulso fica indefinido e o \csname vira \relax.
%
% Motor: o projeto é lualatex (fontspec/polyglossia).  Se o editor chamar pdflatex, cai
% num preâmbulo equivalente com inputenc/babel — serve para conferir, não para o PDF final.
\documentclass[11pt,a5paper]{book}
\ifdefined\directlua
  \usepackage{fontspec}
  \usepackage{polyglossia}
  \setmainlanguage{portuguese}
\else\ifdefined\XeTeXrevision
  \usepackage{fontspec}
  \usepackage{polyglossia}
  \setmainlanguage{portuguese}
\else
  \usepackage[utf8]{inputenc}
  \usepackage[T1]{fontenc}
  \usepackage{lmodern}
  \usepackage[portuguese]{babel}
  \DeclareUnicodeCharacter{2192}{\ensuremath{\rightarrow}}
\fi\fi
% raiz do repo: compilado da raiz ou de volN/<parte>/ (dois níveis abaixo)
\IfFileExists{tex/preambulo.tex}{% Preâmbulo comum ao livro e à pré-visualização de fragmentos (tex/_preview.tex).

% Mancha em A5.  Sem isto vale o padrão da classe book (~105 mm de largura e margem
% externa de ~25 mm em frente e verso), largo demais para o formato.  Aqui: ~118 mm de
% largura, margem interna um pouco maior que a externa (encadernação).
\usepackage[a5paper,inner=16mm,outer=14mm,top=14mm,bottom=16mm,
            includehead,headsep=5mm,footskip=8mm]{geometry}

% Títulos de secção sem justificar (evita linhas estouradas em \Large bold, p. ex.
% «Demonstrabilidade») e um pouco de folga para enunciados de tese em negrito.
\usepackage{sectsty}
\allsectionsfont{\raggedright}
\setlength{\emergencystretch}{1.5em}

% Marca da página impressa do original: \pg{286} imprime [286] no fluxo do texto.
\newcommand{\pg}[1]{\,{\footnotesize\textbf{[#1]}}\,}

% Virada de página do original DENTRO de uma nota (\pgn{422}); discreta, para não quebrar a nota.
\newcommand{\pgn}[1]{{\scriptsize\textbf{[#1]}}}

% Asterisco do original (divisão de assunto, sem título): \asterisco centrado em linha própria,
% também válido no meio do parágrafo ou dentro de nota (\footnote{... \asterisco ...}).
\newcommand{\asterisco}{\par\medskip\centerline{\footnotesize$\ast$}\medskip\noindent}

% Gravura recortada do fac-símile (<figure> do TEI; tools/tei2latex.py): fora de float, no ponto
% exato do original.  \gravura[detalhes]{arquivo.png}{fração da largura do texto}{legenda}:
% legenda principal centrada ("FIGURA VIII. — Côrte…"); os detalhes opcionais (<p> do <figure>:
% "A, hemispherio…; — B, …") vão abaixo, em corpo menor e justificados.  Ambos podem ser vazios.
\usepackage{graphicx}
\newcommand{\gravura}[4][]{\par\medskip\begingroup\centering
  \includegraphics[width=#3\linewidth,height=0.75\textheight,keepaspectratio]{#2}\par
  \if\relax\detokenize{#4}\relax\else\smallskip{\small #4\par}\fi
  \if\relax\detokenize{#1}\relax\else\smallskip{\footnotesize\leftskip=1em\rightskip=1em
    \parfillskip=0pt plus 1fil\noindent #1\par}\fi\endgroup\medskip\noindent}

% Citações latinas de S. Tomás (o OCR erra muito aqui: sempre conferir na imagem).
\newenvironment{latim}{\begin{quote}\itshape}{\end{quote}}

% --- Estrutura do autor (promovido de "Artigo III.tex") -----------------------
% enunciado de tese numerada que corre com o texto (segue-se travessão)
\newcommand{\tese}[2]{\bigskip\noindent\textbf{#1.\quad #2}\;\textemdash\;}
% enunciado de tese numerada em linha própria (seguem-se as subdivisões a), b)...)
\newcommand{\teseB}[2]{\bigskip\noindent\textbf{#1.\quad #2}\par}
% subdivisão a), b), c)... com rubrica em itálico
\newcommand{\sub}[2]{\medskip\noindent#1)\quad\emph{#2},\;\textemdash\;}
% subdivisão a), b), c)... sem rubrica
\newcommand{\subp}[1]{\medskip\noindent#1)\quad}

% --- Notas: no pé (padrão) ou correndo no corpo do texto ---------------------
% Passar  \def\notascorridas{}  no -usepretex faz cada nota correr no corpo, em
% corpo menor, num bloco recuado que interrompe a leitura até acabar.  Ex.:
%   latexmk -lualatex -usepretex='\def\fragmento{vol2/.../01_cap_y.tex}\def\notascorridas{}' tex/_preview.tex
\newif\ifnotascorridas
\ifdefined\notascorridas \notascorridastrue \fi
\ifnotascorridas
  \renewcommand\footnote[1]{%
    \stepcounter{footnote}\textsuperscript{\thefootnote}% marca no ponto do texto
    \par
    {\footnotesize
     \leftskip=2em \rightskip=2em \parindent=1em
     \addvspace{\medskipamount}%
     \noindent\textit{Nota \thefootnote.}\enspace #1\par
     \addvspace{\medskipamount}}%
  }
\fi

% --- Sem notas: \def\semnotas{} no -usepretex suprime todas as notas ---------
% (nem a chamada no texto nem o conteúdo).  Usado por  ./compila.sh --sem-notas .
\ifdefined\semnotas
  \renewcommand\footnote[2][]{}%
  \renewcommand\footnotemark[1][]{}%
  \renewcommand\footnotetext[2][]{}%
\fi

% peças pré-textuais (nó pretextual; p@type em web/obra/esquema.py)
\newcommand{\preTituloObra}[1]{\par\vspace{1em}{\centering\Large\bfseries\scshape #1\par}\vspace{1em}}
\newcommand{\preCentro}[1]{\par{\centering #1\par}}
\newcommand{\preEnderecamento}[1]{\par\vspace{1em}{\centering\small\scshape #1\par}}
\newcommand{\preSaudacao}[1]{\par\medskip\noindent #1\par}
\newcommand{\preDataLocal}[1]{\par\medskip\noindent\emph{#1}\par}
\newcommand{\preFecho}[1]{\par\smallskip{\raggedleft\scshape #1\par}}
}{\input{../../tex/preambulo}}
\def\fimavulso{\end{document}}
\begin{document}
}{\input{../../tex/avulso}}\fi
%
% e termina com  \csname fimavulso\endcsname .  Dentro do livro ou do _preview, \pg já existe
% e nada disto roda; \fimavulso fica indefinido e o \csname vira \relax.
%
% Motor: o projeto é lualatex (fontspec/polyglossia).  Se o editor chamar pdflatex, cai
% num preâmbulo equivalente com inputenc/babel — serve para conferir, não para o PDF final.
\documentclass[11pt,a5paper]{book}
\ifdefined\directlua
  \usepackage{fontspec}
  \usepackage{polyglossia}
  \setmainlanguage{portuguese}
\else\ifdefined\XeTeXrevision
  \usepackage{fontspec}
  \usepackage{polyglossia}
  \setmainlanguage{portuguese}
\else
  \usepackage[utf8]{inputenc}
  \usepackage[T1]{fontenc}
  \usepackage{lmodern}
  \usepackage[portuguese]{babel}
  \DeclareUnicodeCharacter{2192}{\ensuremath{\rightarrow}}
\fi\fi
% raiz do repo: compilado da raiz ou de volN/<parte>/ (dois níveis abaixo)
\IfFileExists{tex/preambulo.tex}{% Preâmbulo comum ao livro e à pré-visualização de fragmentos (tex/_preview.tex).

% Mancha em A5.  Sem isto vale o padrão da classe book (~105 mm de largura e margem
% externa de ~25 mm em frente e verso), largo demais para o formato.  Aqui: ~118 mm de
% largura, margem interna um pouco maior que a externa (encadernação).
\usepackage[a5paper,inner=16mm,outer=14mm,top=14mm,bottom=16mm,
            includehead,headsep=5mm,footskip=8mm]{geometry}

% Títulos de secção sem justificar (evita linhas estouradas em \Large bold, p. ex.
% «Demonstrabilidade») e um pouco de folga para enunciados de tese em negrito.
\usepackage{sectsty}
\allsectionsfont{\raggedright}
\setlength{\emergencystretch}{1.5em}

% Marca da página impressa do original: \pg{286} imprime [286] no fluxo do texto.
\newcommand{\pg}[1]{\,{\footnotesize\textbf{[#1]}}\,}

% Virada de página do original DENTRO de uma nota (\pgn{422}); discreta, para não quebrar a nota.
\newcommand{\pgn}[1]{{\scriptsize\textbf{[#1]}}}

% Asterisco do original (divisão de assunto, sem título): \asterisco centrado em linha própria,
% também válido no meio do parágrafo ou dentro de nota (\footnote{... \asterisco ...}).
\newcommand{\asterisco}{\par\medskip\centerline{\footnotesize$\ast$}\medskip\noindent}

% Gravura recortada do fac-símile (<figure> do TEI; tools/tei2latex.py): fora de float, no ponto
% exato do original.  \gravura[detalhes]{arquivo.png}{fração da largura do texto}{legenda}:
% legenda principal centrada ("FIGURA VIII. — Côrte…"); os detalhes opcionais (<p> do <figure>:
% "A, hemispherio…; — B, …") vão abaixo, em corpo menor e justificados.  Ambos podem ser vazios.
\usepackage{graphicx}
\newcommand{\gravura}[4][]{\par\medskip\begingroup\centering
  \includegraphics[width=#3\linewidth,height=0.75\textheight,keepaspectratio]{#2}\par
  \if\relax\detokenize{#4}\relax\else\smallskip{\small #4\par}\fi
  \if\relax\detokenize{#1}\relax\else\smallskip{\footnotesize\leftskip=1em\rightskip=1em
    \parfillskip=0pt plus 1fil\noindent #1\par}\fi\endgroup\medskip\noindent}

% Citações latinas de S. Tomás (o OCR erra muito aqui: sempre conferir na imagem).
\newenvironment{latim}{\begin{quote}\itshape}{\end{quote}}

% --- Estrutura do autor (promovido de "Artigo III.tex") -----------------------
% enunciado de tese numerada que corre com o texto (segue-se travessão)
\newcommand{\tese}[2]{\bigskip\noindent\textbf{#1.\quad #2}\;\textemdash\;}
% enunciado de tese numerada em linha própria (seguem-se as subdivisões a), b)...)
\newcommand{\teseB}[2]{\bigskip\noindent\textbf{#1.\quad #2}\par}
% subdivisão a), b), c)... com rubrica em itálico
\newcommand{\sub}[2]{\medskip\noindent#1)\quad\emph{#2},\;\textemdash\;}
% subdivisão a), b), c)... sem rubrica
\newcommand{\subp}[1]{\medskip\noindent#1)\quad}

% --- Notas: no pé (padrão) ou correndo no corpo do texto ---------------------
% Passar  \def\notascorridas{}  no -usepretex faz cada nota correr no corpo, em
% corpo menor, num bloco recuado que interrompe a leitura até acabar.  Ex.:
%   latexmk -lualatex -usepretex='\def\fragmento{vol2/.../01_cap_y.tex}\def\notascorridas{}' tex/_preview.tex
\newif\ifnotascorridas
\ifdefined\notascorridas \notascorridastrue \fi
\ifnotascorridas
  \renewcommand\footnote[1]{%
    \stepcounter{footnote}\textsuperscript{\thefootnote}% marca no ponto do texto
    \par
    {\footnotesize
     \leftskip=2em \rightskip=2em \parindent=1em
     \addvspace{\medskipamount}%
     \noindent\textit{Nota \thefootnote.}\enspace #1\par
     \addvspace{\medskipamount}}%
  }
\fi

% --- Sem notas: \def\semnotas{} no -usepretex suprime todas as notas ---------
% (nem a chamada no texto nem o conteúdo).  Usado por  ./compila.sh --sem-notas .
\ifdefined\semnotas
  \renewcommand\footnote[2][]{}%
  \renewcommand\footnotemark[1][]{}%
  \renewcommand\footnotetext[2][]{}%
\fi

% peças pré-textuais (nó pretextual; p@type em web/obra/esquema.py)
\newcommand{\preTituloObra}[1]{\par\vspace{1em}{\centering\Large\bfseries\scshape #1\par}\vspace{1em}}
\newcommand{\preCentro}[1]{\par{\centering #1\par}}
\newcommand{\preEnderecamento}[1]{\par\vspace{1em}{\centering\small\scshape #1\par}}
\newcommand{\preSaudacao}[1]{\par\medskip\noindent #1\par}
\newcommand{\preDataLocal}[1]{\par\medskip\noindent\emph{#1}\par}
\newcommand{\preFecho}[1]{\par\smallskip{\raggedleft\scshape #1\par}}
}{\input{../../tex/preambulo}}
\def\fimavulso{\end{document}}
\begin{document}
}{\input{../../tex/avulso}}\fi
%
% e termina com  \csname fimavulso\endcsname .  Dentro do livro ou do _preview, \pg já existe
% e nada disto roda; \fimavulso fica indefinido e o \csname vira \relax.
%
% Motor: o projeto é lualatex (fontspec/polyglossia).  Se o editor chamar pdflatex, cai
% num preâmbulo equivalente com inputenc/babel — serve para conferir, não para o PDF final.
\documentclass[11pt,a5paper]{book}
\ifdefined\directlua
  \usepackage{fontspec}
  \usepackage{polyglossia}
  \setmainlanguage{portuguese}
\else\ifdefined\XeTeXrevision
  \usepackage{fontspec}
  \usepackage{polyglossia}
  \setmainlanguage{portuguese}
\else
  \usepackage[utf8]{inputenc}
  \usepackage[T1]{fontenc}
  \usepackage{lmodern}
  \usepackage[portuguese]{babel}
  \DeclareUnicodeCharacter{2192}{\ensuremath{\rightarrow}}
\fi\fi
% raiz do repo: compilado da raiz ou de volN/<parte>/ (dois níveis abaixo)
\IfFileExists{tex/preambulo.tex}{% Preâmbulo comum ao livro e à pré-visualização de fragmentos (tex/_preview.tex).

% Mancha em A5.  Sem isto vale o padrão da classe book (~105 mm de largura e margem
% externa de ~25 mm em frente e verso), largo demais para o formato.  Aqui: ~118 mm de
% largura, margem interna um pouco maior que a externa (encadernação).
\usepackage[a5paper,inner=16mm,outer=14mm,top=14mm,bottom=16mm,
            includehead,headsep=5mm,footskip=8mm]{geometry}

% Títulos de secção sem justificar (evita linhas estouradas em \Large bold, p. ex.
% «Demonstrabilidade») e um pouco de folga para enunciados de tese em negrito.
\usepackage{sectsty}
\allsectionsfont{\raggedright}
\setlength{\emergencystretch}{1.5em}

% Marca da página impressa do original: \pg{286} imprime [286] no fluxo do texto.
\newcommand{\pg}[1]{\,{\footnotesize\textbf{[#1]}}\,}

% Virada de página do original DENTRO de uma nota (\pgn{422}); discreta, para não quebrar a nota.
\newcommand{\pgn}[1]{{\scriptsize\textbf{[#1]}}}

% Asterisco do original (divisão de assunto, sem título): \asterisco centrado em linha própria,
% também válido no meio do parágrafo ou dentro de nota (\footnote{... \asterisco ...}).
\newcommand{\asterisco}{\par\medskip\centerline{\footnotesize$\ast$}\medskip\noindent}

% Gravura recortada do fac-símile (<figure> do TEI; tools/tei2latex.py): fora de float, no ponto
% exato do original.  \gravura[detalhes]{arquivo.png}{fração da largura do texto}{legenda}:
% legenda principal centrada ("FIGURA VIII. — Côrte…"); os detalhes opcionais (<p> do <figure>:
% "A, hemispherio…; — B, …") vão abaixo, em corpo menor e justificados.  Ambos podem ser vazios.
\usepackage{graphicx}
\newcommand{\gravura}[4][]{\par\medskip\begingroup\centering
  \includegraphics[width=#3\linewidth,height=0.75\textheight,keepaspectratio]{#2}\par
  \if\relax\detokenize{#4}\relax\else\smallskip{\small #4\par}\fi
  \if\relax\detokenize{#1}\relax\else\smallskip{\footnotesize\leftskip=1em\rightskip=1em
    \parfillskip=0pt plus 1fil\noindent #1\par}\fi\endgroup\medskip\noindent}

% Citações latinas de S. Tomás (o OCR erra muito aqui: sempre conferir na imagem).
\newenvironment{latim}{\begin{quote}\itshape}{\end{quote}}

% --- Estrutura do autor (promovido de "Artigo III.tex") -----------------------
% enunciado de tese numerada que corre com o texto (segue-se travessão)
\newcommand{\tese}[2]{\bigskip\noindent\textbf{#1.\quad #2}\;\textemdash\;}
% enunciado de tese numerada em linha própria (seguem-se as subdivisões a), b)...)
\newcommand{\teseB}[2]{\bigskip\noindent\textbf{#1.\quad #2}\par}
% subdivisão a), b), c)... com rubrica em itálico
\newcommand{\sub}[2]{\medskip\noindent#1)\quad\emph{#2},\;\textemdash\;}
% subdivisão a), b), c)... sem rubrica
\newcommand{\subp}[1]{\medskip\noindent#1)\quad}

% --- Notas: no pé (padrão) ou correndo no corpo do texto ---------------------
% Passar  \def\notascorridas{}  no -usepretex faz cada nota correr no corpo, em
% corpo menor, num bloco recuado que interrompe a leitura até acabar.  Ex.:
%   latexmk -lualatex -usepretex='\def\fragmento{vol2/.../01_cap_y.tex}\def\notascorridas{}' tex/_preview.tex
\newif\ifnotascorridas
\ifdefined\notascorridas \notascorridastrue \fi
\ifnotascorridas
  \renewcommand\footnote[1]{%
    \stepcounter{footnote}\textsuperscript{\thefootnote}% marca no ponto do texto
    \par
    {\footnotesize
     \leftskip=2em \rightskip=2em \parindent=1em
     \addvspace{\medskipamount}%
     \noindent\textit{Nota \thefootnote.}\enspace #1\par
     \addvspace{\medskipamount}}%
  }
\fi

% --- Sem notas: \def\semnotas{} no -usepretex suprime todas as notas ---------
% (nem a chamada no texto nem o conteúdo).  Usado por  ./compila.sh --sem-notas .
\ifdefined\semnotas
  \renewcommand\footnote[2][]{}%
  \renewcommand\footnotemark[1][]{}%
  \renewcommand\footnotetext[2][]{}%
\fi

% peças pré-textuais (nó pretextual; p@type em web/obra/esquema.py)
\newcommand{\preTituloObra}[1]{\par\vspace{1em}{\centering\Large\bfseries\scshape #1\par}\vspace{1em}}
\newcommand{\preCentro}[1]{\par{\centering #1\par}}
\newcommand{\preEnderecamento}[1]{\par\vspace{1em}{\centering\small\scshape #1\par}}
\newcommand{\preSaudacao}[1]{\par\medskip\noindent #1\par}
\newcommand{\preDataLocal}[1]{\par\medskip\noindent\emph{#1}\par}
\newcommand{\preFecho}[1]{\par\smallskip{\raggedleft\scshape #1\par}}
}{\input{../../tex/preambulo}}
\def\fimavulso{\end{document}}
\begin{document}
